\scanpage{101}
\mspage{18}
\margin{De plus, il y a lieu d'exiger \unsure{1°} que $\Theta$ soit trivial sur $\unsure{I'\cap X'_K}$, et qu'une puissance $\Theta^{\unsure{\nu K}}$ de $\Theta$ $(\nu\geq1)$ soit \unsure{triviale} sur $I$. Car $X'_K\ne\varnothing$ implique alors que $X'$ est $\unsure{+}$, $\unsure{=}$ connexe. \deleted{De plus}}

de $\Theta$ se relève en une action de $\Theta$ sur $\widehat{\Pi}_1(X',I')$, où $I'=X'\vert I$.
\deleted{$\unsure{\exists a'\in}$} $I'\cap X'_K\ne\varnothing$, $X'$ \unsure{étant|ét.} connexe (il suffit que $\Gamma$ soit transitif sur $\pi_0(X')$ i.e. que $X'$ soit ``$K$-connexe'' i.e. connexe en tant que schéma sur $K$) alors l'action remontée est unique.

\inserted{La condition 2 s'explicite en pratique comme une condition de stabilité de certains ss-groupes ouverts dans $\pi_1(X,x)$ $(x\in X_K\cap I)$ sous l'action de $\unsure{\pi|\Gamma}$.}

3) Sous les conditions de 2), pour $\forall$ hom $f:X'\longrightarrow Y$ et \deleted{tout}\inserted{un} $x'\in X'_K\cap I'$, \deleted{l'action} \unsure{on} peut trouver un $\Theta'$ \deleted{\unsure{des tels}} de telle façon que l'autom. de $\widehat{\pi}_1(X',x')$ qu'il induit \deleted{\unsure{laisse}} stabilise le noyau de
\[
\widehat{\pi}_1(X',x')\mathbin{\unsure{-|\longrightarrow}}\widehat{\pi}_1(Y,f(x)).
\]
C'est encore une condition de stabilité pour certains sous-groupes, mais \unsure{pas néc.} d'indice fini, de $\pi_1(X,x)$ par $\Theta$.

\margin{On peut \deleted{\unsure{généraliser}}\inserted{\unsure{renforcer}} cette condition, en \unsure{en supposant} \deleted{$\unsure{\forall}$} $x'\in X'_K\cap I'$, et en \deleted{\unsure{l'imposant}} exigeant que pour $\forall\nu$ tel $\unsure{\Theta'^\nu|\Theta^\nu}x'=x'$, $\Theta'^\nu$ sur $\pi_1(X',x')$, \unsure{préserve le noyau de} $\widehat{\pi}_1(\unsure{Y},\unsure{f(x')})$.}

4) Soient $X'_1,X'_2$ deux rev. étales finis de $X$, et $f$ un morphisme
\[
X'_1\xrightarrow{f}X'_2
\]
(\uuline{pas néc. un $X$-morphisme !}), $I'_1=X'_1\vert I$, $I'_2=X'_2\vert I$ on suppose \deleted{\unsure{déjà que}} que $J\overset{\mathrm{def}}=I'_1\cap f^{-1}(I'_2)\ne\varnothing$, \deleted{plus précisément que} \deleted{$I'_1\cap f^{-1}(I'_2)\longrightarrow$} alors \deleted{\unsure{le morph}} le morphisme
\[
\unsure{\Pi_1|\widehat{\Pi}_1}(X'_1,J)\longrightarrow\widehat{\Pi}_1(X'_2,I'_2)
\]
est compatible avec l'action \inserted{prolongée} de $\Theta$ sur les groupoïdes (pour un choix compatible de cette action, s'il y a \unsure{indétermination}, cf 2°).
\mspage{100}
